\documentclass[a4paper,9pt]{extarticle}

%=========================================================
% PACKAGES
%=========================================================

\usepackage[
a4paper,
top=1.8cm,
bottom=1.5cm,
footskip=7mm,
left=1cm,
right=1cm,
headheight=14pt,
headsep=5mm
]{geometry}

\usepackage[T1]{fontenc}
\usepackage[utf8]{inputenc}
\usepackage[french]{babel}

\usepackage[sfdefault]{FiraSans}
\renewcommand{\familydefault}{\sfdefault}

\usepackage{multicol}
\setlength{\columnsep}{0.7cm}

\usepackage{amsmath,amssymb}
\usepackage{titlesec}
\usepackage{enumitem}

\usepackage{tcolorbox}
\tcbuselibrary{skins,breakable}

\usepackage{xcolor}
\usepackage{graphicx}
\usepackage{float}

\usepackage{fancyhdr}
\usepackage[hidelinks]{hyperref}


%=========================================================
% PAGE SETTINGS
%=========================================================

\setlength{\parindent}{0pt}
\setlength{\parskip}{1.5pt}

\setlength{\abovedisplayskip}{4pt}
\setlength{\belowdisplayskip}{4pt}

\setlist[itemize]{
leftmargin=*,
itemsep=1pt,
topsep=1pt
}
\setlist[enumerate]{
leftmargin=*,
itemsep=1pt,
topsep=1pt
}


%=========================================================
% COLORS
%=========================================================

\definecolor{Blue}{HTML}{E3F0FF}
\definecolor{Green}{HTML}{E3F7EA}
\definecolor{Pink}{HTML}{FFE5ED}
\definecolor{Orange}{HTML}{FFF0DC}
\definecolor{Purple}{HTML}{EEE7FF}
\definecolor{Yellow}{HTML}{FFFADB}
\definecolor{Goat}{HTML}{F7EEDF}

\definecolor{BlueTitle}{HTML}{285B8F}
\definecolor{GreenTitle}{HTML}{287052}
\definecolor{PinkTitle}{HTML}{A84063}
\definecolor{OrangeTitle}{HTML}{A96500}
\definecolor{PurpleTitle}{HTML}{6245A5}
\definecolor{GoatTitle}{HTML}{8A5A1F}
\definecolor{RecapTitle}{HTML}{5B3A8E}


%=========================================================
% HEADER
%=========================================================

\pagestyle{fancy}
\fancyhf{}

\fancyhead[L]{\small\textbf{Comportement mécanique des matériaux}}
\fancyhead[C]{\small Présentation, essai de traction, striction}
\fancyhead[R]{\small Week 01 --- 07-13/09/2026}

\fancyfoot[C]{\small\href{https://github.com/AlphaXT38G/notes_ba5}{github.com/AlphaXT38G/notes\_ba5}}
\fancyfoot[R]{\small \thepage}

\renewcommand{\headrulewidth}{0.4pt}


%=========================================================
% TITLES
%=========================================================

\titleformat{\section}
{\large\bfseries\color{BlueTitle}}
{}{0pt}{}

\titlespacing{\section}
{0pt}{6pt}{3pt}


%=========================================================
% BOXES
%=========================================================
% Boxes are kept whole (unbreakable): inside multicols, a box split by tcolorbox
% leaves an empty column behind it.

\newtcolorbox{definition}{
enhanced, unbreakable, colback=Blue, colframe=BlueTitle,
boxrule=0.7pt, arc=0mm, title=\textbf{Définition}, fonttitle=\bfseries
}

\newtcolorbox{important}{
enhanced, unbreakable, colback=Pink, colframe=PinkTitle,
boxrule=0.7pt, arc=0mm, title=\textbf{Idée clé}, fonttitle=\bfseries
}

\newtcolorbox{formula}{
enhanced, unbreakable, colback=Green, colframe=GreenTitle,
boxrule=0.7pt, arc=0mm, title=\textbf{Formules clés}, fonttitle=\bfseries
}

\newtcolorbox{warning}{
enhanced, unbreakable, colback=Orange, colframe=OrangeTitle,
boxrule=0.7pt, arc=0mm, title=\textbf{Attention}, fonttitle=\bfseries
}

\newtcolorbox{example}{
enhanced, unbreakable, colback=Purple, colframe=PurpleTitle,
boxrule=0.7pt, arc=0mm, title=\textbf{Exemple}, fonttitle=\bfseries
}

% GOAT - brown: the idea told as simply as possible
\newtcolorbox{goat}{
enhanced, unbreakable, colback=Goat, colframe=GoatTitle,
boxrule=0.7pt, arc=0mm, title=\textbf{Expliqué à une chèvre}, fonttitle=\bfseries
}

% RECAP - violet: what to know for the exam, at the end of the week
\newtcolorbox{recap}{
enhanced, breakable, colback=Purple, colframe=RecapTitle,
boxrule=1pt, arc=0mm, title=\textbf{Récap pour l'examen}, fonttitle=\bfseries
}

% Screenshot of a slide of the course: \slide{cNN-PP}
\newcommand{\slide}[2][0.95]{%
\begin{center}
\includegraphics[width=#1\linewidth]{../images/course/#2}
\end{center}}


%=========================================================
% DOCUMENT
%=========================================================

\begin{document}

\thispagestyle{fancy}

\begin{multicols}{2}

%=========================================================
\section{Le cours en bref}

\begin{definition}
\textbf{MSE 234} : comportement mécanique, élaboration, structure et cycle de vie des grandes
classes de matériaux de structure : essais mécaniques, céramiques, métaux (Al, Cu, Fe), polymères,
composites et matériaux poreux.
\end{definition}

\begin{warning}
\textbf{Examen sans document ni formulaire} : la liste des formules à connaître par c\oe ur est
sur Moodle. Note : 75\,\% cours + 18\,\% questions de TP (examen) + 7\,\% cahier de labo et
participation aux 7 TP. Ce qui fait foi : les diapositives.
\end{warning}

\begin{goat}
Un cours pour savoir pourquoi un verre casse, pourquoi un trombone plie, et comment on rend un
métal plus solide. À l'examen, rien sous les yeux : les formules doivent être dans la tête.
\end{goat}

%=========================================================
\section{L'essai de traction}

\slide[0.8]{c01-06}

\begin{important}
Dans la section réduite : \textbf{une seule contrainte}, $\sigma_{xx}$, uniforme. On sait
exactement ce qu'on impose. La déformation, elle, a 3 composantes (le barreau maigrit).
\end{important}

\begin{formula}
\[ \varepsilon_{xx} = \frac{\sigma_{xx}}{E}, \qquad \varepsilon_{yy} = \varepsilon_{zz} = -\nu\frac{\sigma_{xx}}{E} \]
Volume constant (plasticité) : $\frac{\Delta V}{V} \approx \varepsilon_{xx}+\varepsilon_{yy}+\varepsilon_{zz} = 0$
$\Rightarrow \varepsilon_{yy} = \varepsilon_{zz} = -\varepsilon_{xx}/2$ ($\nu = 1/2$).
\end{formula}

Mesures : force par cellule de charge ; allongement par traverse (imprécis), extensomètre, jauge
ou suivi optique.

\begin{goat}
On tire sur un barreau aminci au milieu : là, toute la matière est tirée pareil, dans une seule
direction. Il s'allonge et maigrit, comme un chewing-gum.
\end{goat}

%=========================================================
\section{Nominal vs vrai}

\begin{formula}
\[ e = \frac{L-L_0}{L_0}, \quad \varepsilon = \int_{L_0}^{L}\frac{dL}{L} = \ln\frac{L}{L_0} = \ln(1+e) \]
\[ R = \frac{P}{S_0}, \quad \sigma = \frac{P}{S} \approx \frac{P}{S_0}\frac{L}{L_0} = R(1+e) \]
(la seconde : volume constant $S_0L_0 = SL$).
\end{formula}

\begin{important}
Les déformations vraies s'additionnent ; $\varepsilon \approx e$ aux petites déformations.
$\sigma > R$ car la section diminue.
\end{important}

\begin{example}
\textit{(pas dans les diapos)} $L_0 = 50 \to L = 60$\,mm, $P = 20$\,kN, $S_0 = 50$\,mm$^2$ :
$e = 0{,}20$, $\varepsilon = 0{,}182$, $R = 400$\,MPa, $\sigma = 480$\,MPa.
\end{example}

\begin{goat}
Nominal : on fait comme si le barreau n'avait pas bougé. Vrai : on le regarde tel qu'il est,
plus long et plus maigre, donc plus tiré pour la même force.
\end{goat}

%=========================================================
\section{Lire la courbe conventionnelle}

\slide[0.85]{c01-14}

\begin{definition}
$E$ (pente) ; $R_e$ (premier écart, mal défini) ; $\mathbf{R_{0,2}}$ (parallèle décalée de
0,2\,\%) ; $\mathbf{R_m}$ (maximum, UTS) ; $e_h$ (déformation répartie) ; $e_F$ (après rupture).
\end{definition}

\slide[0.85]{c01-15}

Allures : (a) fragile, (b) métal ductile, (c) crochet + bandes de Lüders (acier doux,
polymères), (d) élastomère (500--1000\,\%, réversible).

\begin{important}
\begin{itemize}
\item Décharge : droite de pente $E$ ; seule la part élastique revient.
\item Énergie volumique $\mathcal{E} = \int_0^\varepsilon\sigma\,d\varepsilon$ (J/m$^3$)
      = élastique (rendue) + irréversible (chaleur, dislocations).
\item Contrainte d'écoulement $\sigma(\varepsilon)$ ; taux d'écrouissage
      $\theta = d\sigma/d\varepsilon$.
\item Ductilité : $\Delta L/L_0 \le 1\,\%$ fragile, $> 10\,\%$ ductile ; ou $A_f/A_0$.
\end{itemize}
\end{important}

\begin{goat}
D'abord un ressort, puis le barreau s'étire pour de bon en durcissant, puis il s'amincit à un
endroit et casse. Si on lâche en route, seul le \og ressort\fg{} revient.
\end{goat}

\end{multicols}

\end{document}
